% Quarantine-only source-aware proof.
% Date: 29 August 2026.
% This file is not a live-book input and makes no novelty or priority claim.

\begin{proof}[Proof of
\cref{lem:p21p3-componentwise-normal-rigid-zero-atom}]
We first reduce normality to the precise local coordinate needed in the
estimate.  Since $X$ is good, $\mathcal O_{X,y}$ is a noetherian henselian
excellent local ring by \cite[\S11.1.10]{CLD25}.  A rigid point is Zariski
closed in a sufficiently small strict affinoid neighborhood.  Indeed, the
fraction field of the quotient by the evaluation kernel embeds densely in
the finite-dimensional $K$-vector space $\mathscr H(y)$.  It is therefore
finite-dimensional over $K$, as is its $K$-subalgebra given by the affinoid
quotient.  That quotient is a finite-dimensional domain over $K$, hence a
field, so the evaluation kernel is maximal.  The central dimension of the
germ is consequently zero, and the local dimension formula
\cite[(11.5.5.1)]{CLD25}, together with the assumed pure local dimension,
gives
\[
                  \dim_{\mathrm{Krull}}\mathcal O_{X,y}=1.
\]
The normal local ring $\mathcal O_{X,y}$ is a domain.  The equivalence
between a one-dimensional noetherian normal local domain and a discrete
valuation ring is
\cite[\href{https://stacks.math.columbia.edu/tag/00PD}{Tag~00PD}]
{stacks-project}.  We may therefore choose a uniformizer
$\pi\in\mathcal O_{X,y}$.

Put $X^\circ:=\operatorname{Int}(X/K)$ and
$\Omega^\circ:=\Omega\cap X^\circ$.  These sets are open, and the
boundary-composition identity \cite[(11.4.1.3)]{CLD25} shows that every open
analytic subspace of $X^\circ$ is boundaryless.  After shrinking around
$y$, we may consequently work in a boundaryless good separated strict open
set
\[
                         U\subseteq\Omega^\circ
\]
on which $X$ is purely one-dimensional and normal and on which $\pi$ is
represented by an analytic function $f$ satisfying $f(y)=0$.  Such a normal
shrinking exists because the ordinary normal locus is open by
\cite[\S11.3.1]{CLD25}.  The germ of the fibre of
$f:U\to\mathbf A_K^{1,\mathrm{an}}$ over $0$ is zero-dimensional at $y$,
because
\[
                  \mathcal O_{X,y}/(f)
                  =\mathcal O_{X,y}/(\pi)
\]
is a field.  Formula (11.4.1.3) also gives
$y\in\operatorname{Int}(U/\mathbf A_K^{1,\mathrm{an}})$.  Thus
\cite[Lemma~11.6.5(a)]{CLD25} makes $f$ finite at $y$.

We spell out the flatness point, since ordinary module-flatness must not be
silently substituted for analytic flatness.  Let
\[
       A:=\mathcal O_{\mathbf A_K^{1,\mathrm{an}},0},
       \qquad B:=\mathcal O_{X,y}.
\]
The ring $A$ is a discrete valuation ring with uniformizer $T$, and the
local homomorphism $A\to B$ sends $T$ to $\pi$.  It is injective: if
$0\ne a\in A$, write $a=T^m v$ with $v$ a unit; then
$a(f)=\pi^m v(f)\ne0$ in the domain $B$.  Since $f$ is finite at $y$, the
$A$-module $B$ is finite by \cite[\S11.4.4]{CLD25}.  It is torsion-free,
hence finite free over the discrete valuation ring $A$.  Thus $B$ is flat
over $A$.  Since $y\in\operatorname{Int}(U/\mathbf A_K^{1,\mathrm{an}})$,
the local-ring criterion of \cite[Remark~11.4.7]{CLD25} upgrades this
module-flatness to analytic flatness of $f$ at $y$.  The local finite-flat
construction of \cite[\S11.6.3]{CLD25} now gives compact analytic
neighborhoods on which $f$ is finite and flat and on which $y$ is the only
point over $0$.  Restricting the target to a sufficiently small closed disc
of admissible radius $R_0>0$, and the source to its inverse image, therefore
produces a compact analytic neighborhood $Q_{R_0}\subset U$ of $y$ and a
finite flat map of constant degree $d>0$
\begin{equation}
\label{eq:p213-normal-rigid-finite-tube}
          f:Q_{R_0}\longrightarrow D(0,R_0),
          \qquad f^{-1}(0)=\{y\}.
\end{equation}
A finite morphism has empty relative boundary.  By
\cite[(11.4.1.3) and Example~11.4.2]{CLD25}, the boundary of $Q_{R_0}$ is
therefore contained in the inverse image of the outer Gauss point
$\{|T|=R_0\}$.

Choose an admissible radius $0<R<R_0$ and a number
\[
             0<\varepsilon_0<\tfrac12\log(R_0/R).
\]
For every admissible $0<r<R$ such that
$\log(R/r)>4\varepsilon_0$, set
\[
                    s_R=-\log R,
                    \qquad s_r=-\log r.
\]
Choose a function $\rho_r\in C^\infty(\mathbf R,[0,1])$ which vanishes on
\[
 (-\infty,s_R-\varepsilon_0]\cup
 [s_r+\varepsilon_0,\infty),
\]
equals $1$ on $[s_R,s_r]$, is nondecreasing on
$[s_R-\varepsilon_0,s_R]$, and is nonincreasing on
$[s_r,s_r+\varepsilon_0]$.  Define
\[
 q_r(t):=\int_{-\infty}^{t}\rho_r(\tau)\,\mathrm d\tau,
 \qquad
 H_r:=\int_{\mathbf R}\rho_r(\tau)\,\mathrm d\tau.
\]
Then
\begin{equation}
\label{eq:p213-normal-rigid-radial-data}
 H_r\geq\log(R/r),
 \qquad
 \int_{\mathbf R}|q_r''(t)|\,\mathrm dt=2.
\end{equation}

On $Q_{R_0}\setminus\{y\}$ put
\[
                         \chi_r=q_r(-\log|f|).
\]
The function $q_r$ is zero on an outer collar of $Q_{R_0}$ and is constant
equal to $H_r$ for sufficiently large arguments.  We may therefore extend
$\chi_r$ by $H_r$ across $y$ and by zero to $U\setminus Q_{R_0}$.  The
extensions are made where the function is constant; on the remaining locus
$f$ is invertible.  The moment description of smooth functions in
\cite[Corollary~15.1.2]{CLD25} therefore shows that $\chi_r$ is a
nonnegative compactly supported CLD-smooth function on $U$.  Moreover,
\begin{equation}
\label{eq:p213-normal-rigid-value-at-y}
                         \chi_r(y)=H_r.
\end{equation}

Let $s=-\log|f|$.  On the two transition tubes one has
\begin{equation}
\label{eq:p213-normal-rigid-radial-dd}
       \mathrm d'\mathrm d''\chi_r
       =q_r''(s)\,\mathrm d's\wedge\mathrm d''s.
\end{equation}
The two tubes are disjoint, and $q_r''$ has a fixed sign on each one.
Apply the calibrated pushforward formula
\cite[Proposition~14.4.2]{CLD25} to the restriction of
\eqref{eq:p213-normal-rigid-finite-tube} over each corresponding target
annulus.  The identification of calibrated degree with finite-flat degree
in \cite[Lemma~13.3.8]{CLD25} shows that the pushforward factor is $d$;
this includes all local and residue-degree multiplicities.  Consequently
\begin{align}
 \left\|\mathrm d'\mathrm d''\chi_r\wedge\delta_U
 \right\|_{\mathrm{TV}}
 &=d\int_{\mathbf R}|q_r''(t)|\,\mathrm dt \notag\\
 &=2d.
\label{eq:p213-normal-rigid-radial-tv}
\end{align}
Neither \cite[Proposition~14.4.2]{CLD25} nor
\cite[Lemma~13.3.8]{CLD25} requires the source to be rig-smooth.

Because $\Omega^\circ$ is boundaryless, its integration current is closed
by \cite[\S7.7.4]{CLD25}.  Applying \cite[Lemma~8.5.5]{CLD25} with
$T=\delta_{\Omega^\circ}$ shows that the current
\[
                         \mu_u
       =\mathrm d'\mathrm d''[u|_{\Omega^\circ}]
\]
is positive on $\Omega^\circ$; as a positive current of bidimension
$(0,0)$ it is a Radon measure by \cite[Proposition~7.2.8]{CLD25}.  Its
restriction to the boundaryless open set $U$ is
\[
                 \mu_u|_U=\mathrm d'\mathrm d''[u|_U].
\]
Equivalently, one may apply the same two results with $T=\delta_U$
directly.  Positivity,
\eqref{eq:p213-normal-rigid-radial-data},
\eqref{eq:p213-normal-rigid-value-at-y}, the transposition identity
\cite[Lemma~7.7.5]{CLD25}, and
\eqref{eq:p213-normal-rigid-radial-tv} yield
\begin{align}
 \log(R/r)\,\mu_u(\{y\})
 &\leq H_r\mu_u(\{y\})                                      \notag\\
 &\leq\int_U\chi_r\,\mathrm d\mu_u                          \notag\\
 &=\bigl\langle[u|_U],\mathrm d'\mathrm d''\chi_r\bigr\rangle \notag\\
 &\leq 2d\sup_{Q_{R_0}}|u|.
\label{eq:p213-normal-rigid-final-bound}
\end{align}
The last line follows by first taking the absolute value of the pairing and
then using its total-variation bound.  Its right-hand side is independent of
$r$.  Since the valuation of $K$ is nontrivial, admissible positive radii
tend to zero.  Letting $r\to0$ in
\eqref{eq:p213-normal-rigid-final-bound} proves
\eqref{eq:p213-normal-rigid-zero-atom}.
\end{proof}

% Quarantine-only proof.
% Date: 29 August 2026.
% This file is not a live-book input and makes no novelty or priority claim.

\begin{proof}[Proof of
\cref{lem:p213-trivial-gauss-rigid-zero-atom}]
The normal locus is Zariski open by \cite[\S11.3.1]{CLD25}.  After
shrinking $\Omega$ about $y$, we may therefore assume that
\[
 \Omega\subseteq\operatorname{Int}(X/K),
 \qquad
 \Omega\ \text{is normal and purely one-dimensional}.
\]
An open analytic subspace of the analytic interior is boundaryless by
the boundary-composition formula \cite[(11.4.1.3)]{CLD25}.

Fix $0<r<1$.  Let
\[
 K_r=K((\varpi)),\qquad |\varpi|=r,
\]
with its $\varpi$-adic Gauss norm.  In the notation of
\cite[Example~11.1.3]{CLD25}, this is the field attached to the one-term
polyradius $(r)$.  Since $K$ is trivially valued and $r\ne1$, this
polyradius is $K$-free.  Thus $K_r/K$ is analytically separable by
\cite[\S11.1.9]{CLD25}.  It is also universally multiplicative, and hence
the projection
\[
 \pi:\Omega_r:=\Omega\widehat\otimes_K K_r\longrightarrow\Omega
\]
has the continuous Shilov section
$\sigma:\Omega\to\Omega_r$ described in
\cite[\S12.7.4--Corollary~12.7.5]{CLD25}.

We record all geometric properties needed below.  Analytic separability,
the scalar-extension statements in \cite[\S\S11.1.9 and 11.3.1]{CLD25},
and the $R_1+S_2$ criterion show that $\Omega_r$ is normal.  Every
irreducible component after scalar extension has the same dimension as,
and maps onto, an irreducible component before scalar extension by
\cite[\S11.2.5]{CLD25}; hence $\Omega_r$ is purely one-dimensional.
Boundarylessness is stable under analytic base change by
\cite[\S11.4.1]{CLD25}.  Goodness and separatedness are preserved by
scalar extension.  Strictness follows directly on affinoid charts: the
base change of a quotient of
$K\langle T_1,\ldots,T_m\rangle$ is a quotient of
$K_r\langle T_1,\ldots,T_m\rangle$.

The lift of $y$ is rigid without any separability hypothesis on its
residue field.  Indeed, put $L=\mathscr H(y)$.  Rigidity says that $L/K$
is finite, and the valuation on $L$ is trivial.  Coefficientwise
expansion gives an isometric isomorphism of Banach $K_r$-algebras
\begin{equation}
\label{eq:p213-trivial-gauss-residue-tensor}
             L\widehat\otimes_K K_r\simeq L((\varpi)).
\end{equation}
To see the isometry explicitly, choose a $K$-basis of $L$.  It is
orthonormal for the trivial norm, and the Gauss norm of a Laurent series
with coefficients in $L$ is the maximum of the Gauss norms of its finitely
many coordinate series.  The right-hand side of
\eqref{eq:p213-trivial-gauss-residue-tensor} is a field finite over $K_r$.
Consequently the scalar fibre above $y$ is a singleton, namely
$\sigma(y)$, and
\[
             \mathscr H(\sigma(y))\simeq L((\varpi)).
\]
Thus $\sigma(y)$ is a rigid interior point of the good separated strict
normal pure one-dimensional $K_r$-analytic space $\Omega_r$.  This
argument is unchanged when $L/K$ is purely inseparable.

Set $u_r=\pi^*u$.  It is continuous and CLD-psh by
\cite[Proposition~8.1.5]{CLD25}.  Put
\[
                         \mu_r=\mathrm d'\mathrm d''[u_r].
\]
The scalar-extension support statement
\cite[Proposition~16.8.2]{CLD25}, applied with $n=p=1$ and the smooth
$(0,0)$-form $1$, gives
\begin{equation}
\label{eq:p213-trivial-gauss-shilov-support}
                         \operatorname{supp}(\mu_r)
                         \subseteq\sigma(\Omega).
\end{equation}
The pushforward identity recalled in the same proposition from
\cite[Proposition~8.7.5]{CLD25} gives
\begin{equation}
\label{eq:p213-trivial-gauss-pushforward}
                         \pi_*\mu_r=\mu_u.
\end{equation}
Both are positive Radon measures; positivity and the Radon assertion also
follow directly from \cite[Lemma~8.5.5 and Proposition~7.2.8]{CLD25}.
Since $\pi\circ\sigma$ is the identity,
\eqref{eq:p213-trivial-gauss-shilov-support} and
\eqref{eq:p213-trivial-gauss-pushforward} imply
\begin{equation}
\label{eq:p213-trivial-gauss-atom-transfer}
             \mu_u(\{y\})
             =\mu_r\bigl(\pi^{-1}(\{y\})\bigr)
             =\mu_r(\{\sigma(y)\}).
\end{equation}
The nontrivially valued rigid-point estimate
\cref{lem:p21p3-componentwise-normal-rigid-zero-atom} applies to
$u_r$ at $\sigma(y)$ and makes the last term zero.  This proves
\eqref{eq:p213-trivial-gauss-rigid-zero-atom}.
\end{proof}

% Quarantine-only third-revision candidate proof.
% Date: 29 August 2026.
% This file is not a live-book input and has not been Rethlas-verified.
% It makes no novelty or priority claim.  The rejected first proof and the
% rejected earlier revisions are preserved under their original filenames.

\paragraph{Slope and orientation conventions.}
We use the type decomposition, branch space, residual curve, and canonical
graph metric of \cite[\S\S3.3--4.5]{Ducros24CurveStructure}.  The point
that was mishandled in the rejected first proof is the distinction between
two slopes.  If an analytic unit $f$ maps a sufficiently short section of
a branch $b$ by a dilation of canonical graph metrics with signed ratio
$\varepsilon r$, then
\begin{equation}
             \mathrm d\log_b|f|
                  =\varepsilon r [s(b):K]\in\mathbf Z
\label{eq:p213-normalized-versus-metric-slope}
\end{equation}
by \cite[\S4.5.21.4]{Ducros24CurveStructure}.  Thus the metric derivative
is $\varepsilon r$, not the integer in
\eqref{eq:p213-normalized-versus-metric-slope}.  All lattice and balancing
arguments below use the integer $\mathrm d\log_b|f|$; the canonical metric
is used only to identify the orientation of the image ray.

At a type--$2$ point $y$, the residual curve $C_y$ is projective, normal,
and integral over the actual degree-one residue field of $K$; it need not
be geometrically integral.  Since $y$ is interior, its branches are the
closed points of $C_y$ by
\cite[Theorem~4.2.10 and \S4.2.11.1]{Ducros24CurveStructure}.  A type--$3$
interior point has exactly two branches by
\cite[\S4.2.11.2]{Ducros24CurveStructure}.  For every analytic unit, the
normalized node relation
\begin{equation}
                    \sum_{b}\mathrm d\log_b|f|=0
\label{eq:p213-normalized-node-relation}
\end{equation}
holds on a boundaryless curve by
\cite[Theorem~4.5.22]{Ducros24CurveStructure}.

\begin{proof}[Proof of
\cref{lem:p213-variation-graph-source-cutoff}]
We first separate the active branches by normalized integer slope vectors.
Suppose that $y$ is of type~$2$, write $C=C_y$, and let
\[
                     b_1,\ldots,b_s
\]
be the branches on which $F_0$ is not locally constant.  There are only
finitely many of them: the variation graph of each coordinate of $F_0$ is
locally finite by \cite[\S\S4.5.21.1--4.5.21.2]
{Ducros24CurveStructure}.  Denote by $P_i\in C$ the closed point
corresponding to $b_i$.  If $s=0$, take $S$ to be empty and pass directly
to the source-star and cutoff construction below.  We henceforth assume
$s>0$ in the separation construction.

Weak approximation on $C$ gives rational functions
$r_1,\ldots,r_s\in\operatorname{Rat}(C)^\times$ such that
\begin{equation}
       \operatorname{ord}_{P_j}(r_i)=\delta_{ij}
       \qquad(1\leq i,j\leq s).
\label{eq:p213-variation-first-order-block-revised}
\end{equation}
For completeness, choose an effective divisor $D$ disjoint from all
$P_j$ and so positive that
\[
 H^1\!\left(C,
   \mathcal O_C\!\left(D-2P_i-\sum_{j\ne i}P_j\right)\right)=0.
\]
The resulting jet-evaluation map is surjective.  Prescribing a
uniformizer modulo $\mathfrak m_{P_i}^2$ at $P_i$ and the value $1$
modulo $\mathfrak m_{P_j}$ at every $P_j\ne P_i$ gives $r_i$.
This argument is over the actual constant field of $C$ and requires
neither geometric integrality nor a rational point.

Let $B$ be the finite set of closed points outside
$\{P_1,\ldots,P_s\}$ at which at least one $r_i$ has a zero or a pole.
A second jet-evaluation argument, with a fresh sufficiently positive pole
divisor disjoint from $B\cup\{P_1,\ldots,P_s\}$, gives
$h\in\operatorname{Rat}(C)^\times$ such that
\begin{equation}
 \operatorname{ord}_{P_i}(h)=0,
 \qquad
 \operatorname{ord}_{Q}(h)=1\quad(Q\in B).
\label{eq:p213-variation-collision-killer-revised}
\end{equation}
At any further zero or pole $R$ of $h$, all the $r_i$ have order zero by
the definition of $B$.

We next lift these rational functions and compute their correct slopes.
The field $\operatorname{Rat}(C)$ is the degree-one residue field of
$\mathcal H(y)$.  A norm-one representative in $\mathcal H(y)$ can be
approximated to error strictly less than one by an element of the dense
subfield $\kappa(y)$, without changing its residue.  Since $Y$ is normal
and $y$ is nonrigid,
\begin{equation}
                         \mathcal O_{Y,y}=\kappa(y)
\label{eq:p213-variation-normal-local-field-revised}
\end{equation}
by \cite[\S4.2.18]{Ducros24CurveStructure}.  Hence the approximants are
analytic unit germs.  Shrinking once, all of them are analytic units on a
common neighborhood of $y$.

For a branch $b_P$ corresponding to a closed point $P\in C$, choose a
norm-one lift $s_P$ of a residual uniformizer and orient $b_P$ away from
$y$.  Define
\begin{equation}
             d_P:=-\mathrm d\log_{b_P}|s_P|\in\mathbf Z_{>0}.
\label{eq:p213-variation-integral-branch-weight}
\end{equation}
The positivity follows from the fact that the residue of $s_P$ vanishes
at $P$, so $|s_P|$ decreases on a sufficiently short section of $b_P$.
The number $d_P$ is independent of the chosen norm-one uniformizer.  In
fact, \cite[\S4.2.11.1]{Ducros24CurveStructure} says that the refined
branch valuation is the composite of the original absolute value and the
discrete valuation at $P$.  If $q$ has secondary valuation zero, applying
the inequalities of \cite[\S4.2.6]{Ducros24CurveStructure} to $q$ and
$q^{-1}$ shows that $|q|$ is constant on a short branch section.  Therefore,
for every $g\in\operatorname{Rat}(C)^\times$ and every norm-one lift
$\widehat g$,
\begin{equation}
       -\mathrm d\log_{b_P}|\widehat g|
                 =d_P\operatorname{ord}_P(g).
\label{eq:p213-variation-residual-normalized-slope}
\end{equation}
The same argument for an arbitrary analytic unit germ $f$ gives
\[
                  -\mathrm d\log_{b_P}|f|\in d_P\mathbf Z.
\]

For a finite moment $G=(g_1,\ldots,g_m)$ that is nonconstant on $b$, set
\begin{equation}
 n_b(G):=\bigl(-\mathrm d\log_b|g_1|,\ldots,
                    -\mathrm d\log_b|g_m|\bigr)\in\mathbf Z^m,
 \qquad q_b(G):=\gcd_j|n_{b,j}(G)|.
\label{eq:p213-integral-slope-vector}
\end{equation}
By \eqref{eq:p213-normalized-versus-metric-slope}, the derivative of
$G_{\mathrm{trop}}$ with respect to canonical source length is
$n_b(G)/[s(b):K]$.  Consequently the primitive integral outgoing tangent
of the target ray is
\begin{equation}
                              u_b(G)=n_b(G)/q_b(G).
\label{eq:p213-primitive-integral-tangent}
\end{equation}
This is precisely where the factor $[s(b):K]$ cancels from the direction
but not from the metric dilation.

Let $S$ be the tuple of the lifts of $r_1,\ldots,r_s,h$.  On $b_i$, every
coordinate of $(F_0,S)$ has normalized slope divisible by $d_{P_i}$ and
the $r_i$-coordinate has slope exactly $d_{P_i}$.  Hence
$q_{b_i}(F_0,S)=d_{P_i}$ and the $S$-block of the primitive tangent is
$(e_i,0)$.  At a point $Q\in B$, the last coordinate in
\eqref{eq:p213-variation-collision-killer-revised} is one; at a further
zero or pole of $h$, the first $s$ coordinates vanish and the last does
not.  A branch activated only by $F_0$ has zero $S$-block.  It follows
that each selected ray is distinct from every other source ray.

If $y$ is of type~$3$, choose an element of $\mathcal H(y)^\times$ whose
modulus has non-torsion class modulo $|K^\times|$, approximate it by an
element of $\kappa(y)^\times$ without changing that modulus class, and use
\eqref{eq:p213-variation-normal-local-field-revised} to regard the result
as an analytic unit germ.  The description in
\cite[\S4.2.11.2]{Ducros24CurveStructure} shows that its modulus is
nonconstant on both branches, and
\eqref{eq:p213-normalized-node-relation} says that its two normalized
slopes are opposite.  Appending this one germ separates the two ray
germs.  We again call the appended block $S$.

Put $F=(F_0,S)$.  The union $\Gamma_F$ of the variation graphs of its
coordinates is closed and locally finite, every coordinate is locally
constant off $\Gamma_F$, and $F_{\mathrm{trop}}$ is piecewise a dilation
on $\Gamma_F$ by \cite[\S\S4.5.21.1--4.5.21.2]
{Ducros24CurveStructure}.  Shrink to an open tree $W$ about $y$ such that
$\Gamma_F\cap W$ is a finite star with its only vertex at $y$.  Thus every
active arm is mapped injectively and integral-affinely to a ray, whereas
$F_{\mathrm{trop}}$ is locally constant on every component of
$W\setminus\Gamma_F$.

The space $Y$ is locally holomorphically separated because it is good and
separated.  By \cite[Corollary~15.2.5]{CLD25}, there is a smooth function
$\psi$ with compact support in $W$ that is identically one near $y$.
Choose $\rho\in C^\infty(\mathbf R,[0,1])$ that is zero on a neighborhood
of $0$ and one on a neighborhood of $1$, and put $\chi=\rho\circ\psi$.
This is a smooth function: in every moment chart presenting $\psi$, it is
presented by the composite of the presenting smooth real function with
$\rho$.  Moreover,
\[
 0\leq\chi\leq1,\qquad
 \operatorname{supp}(\chi)\subseteq\operatorname{supp}(\psi)\Subset W,
 \qquad \chi=1\text{ near }y.
\]
Let
\[
        D:=\operatorname{supp}(\mathrm d'\chi)
             \cup\operatorname{supp}(\mathrm d''\chi).
\]
The compact set $D\cap\Gamma_F$ does not contain $y$.  Since every active
arm is mapped injectively away from $v$, one has
\[
                v\notin F_{\mathrm{trop}}(D\cap\Gamma_F).
\]
Choose an open neighborhood $O$ of $v$, small enough that it misses this
compact image and that
$F_{\mathrm{trop}}(\Gamma_F\cap\operatorname{supp}\chi)\cap O$ is the
finite separated fan germ described above.

Let now $T$ be as in the statement.  The support of $\chi T$ is compact,
so \cite[\S7.1.11]{CLD25} defines
$R_{T,F,\chi}=(F_{\mathrm{trop}})_*(\chi T)$ without any properness
assumption on $F_{\mathrm{trop}}$.  It is positive because $\chi\geq0$.
We prove directly that it is closed on $O$.  The Leibniz rule and the
closedness of $T$ give
\begin{equation}
 \mathrm d'R_{T,F,\chi}
     =(F_{\mathrm{trop}})_*(\mathrm d'\chi\wedge T),
 \qquad
 \mathrm d''R_{T,F,\chi}
     =(F_{\mathrm{trop}})_*(\mathrm d''\chi\wedge T).
\label{eq:p213-cutoff-differential}
\end{equation}
To pair either right-hand side with a target test form supported in $O$,
pull that form back by $F_{\mathrm{trop}}$.  On
$D\cap\Gamma_F$ its pullback has disjoint support by the choice of $O$.
On $D\setminus\Gamma_F$ the map $F_{\mathrm{trop}}$ is locally constant,
so the pullback of every positive-degree target form vanishes.  Both
right-hand sides of \eqref{eq:p213-cutoff-differential} therefore vanish
on $O$.

The same local-constancy observation shows that the support of the
pushforward in $O$ is contained in the displayed finite fan: a target
form supported away from the image of $\Gamma_F$ pulls back to zero.
Finally, the primitive $S$-blocks computed above show that a selected ray
germ has only its named source arm above it inside
$\Gamma_F\cap W$.  Its edge coefficient is therefore unmixed.  This
proves the lemma without a global target tube.
\end{proof}

\begin{proof}[Proof of
\cref{lem:p213-euclidean-one-dimensional-current-revised}]
The assertion is local on $\Omega$.  On a relatively compact coordinate
ball meeting only finitely many cells, positivity gives order zero by
\cite[Proposition~7.2.9]{CLD25}.  Thus $R$ is represented by a symmetric
positive-semidefinite matrix $M=(\mu_{ij})$ of Radon measures.  The
equation $\mathrm d'R=0$ says, column by column,
\begin{equation}
                         \sum_i\partial_i\mu_{ij}=0
                 \quad\text{in }\mathcal D'(\Omega).
\label{eq:p213-euclidean-column-divergence-revised}
\end{equation}

Consider a tube meeting the relative interior of one edge.  Choose linear
coordinates $(t,z_2,\ldots,z_N)$ in which the edge is $\{z=0\}$.  In
\eqref{eq:p213-euclidean-column-divergence-revised}, use a scalar test
function $z_a\psi(t,z)$, with arbitrary compactly supported $\psi$, for
each $a\geq2$.  Its tangential derivative vanishes on the edge and its
$a$-th normal derivative there is arbitrary.  Hence
$\mu_{aj}=0$ for $a\geq2$ and all $j$; symmetry gives
$\mu_{ia}=0$ as well.  The remaining equation is
$\partial_t\mu_{11}=0$.  A Radon measure on an interval with zero
distributional derivative is a constant multiple of Lebesgue measure,
and positivity makes the constant nonnegative.  In the integral affine
parametrization by the primitive tangent this is $b_e[e]$.

At a vertex $v$, subtract the edge currents on a finite star.  The
remainder is $A_v\delta_v$ for a positive-semidefinite matrix $A_v$.
The distributional column divergence of the full current at $v$ is, up to
one common sign,
\begin{equation}
 \left(\sum_{e\ni v}b_eu_{e,v}\right)\delta_v
       +\left(\sum_i(A_v)_{ij}\partial_i\delta_v\right)_{j=1}^N.
\label{eq:p213-euclidean-vertex-distribution-revised}
\end{equation}
Smooth test functions have independently prescribed value and gradient at
$v$.  First taking value zero and arbitrary gradient gives $A_v=0$;
then taking arbitrary value gives
\eqref{eq:p213-euclidean-vertex-balance-revised}.  Local finiteness gives
\eqref{eq:p213-euclidean-edge-decomposition-revised}.
\end{proof}

\begin{proof}[Proof of
\cref{prop:p213-normal-curve-arbitrary-valuation-current-revised}]
We first attach an intrinsic ratio to a source branch.  Let $y$ be of
type~$2$ or $3$, let $b$ be a branch at $y$, and choose a finite moment
$F$ supplied by
\cref{lem:p213-variation-graph-source-cutoff} that is nonconstant on $b$
and separates its ray germ.  Write
$n_b(F)$, $q_b(F)$, and $u_b(F)$ as in
\eqref{eq:p213-integral-slope-vector}--
\eqref{eq:p213-primitive-integral-tangent}.  By
\cref{lem:p213-euclidean-one-dimensional-current-revised}, the coefficient
of $(F_{\mathrm{trop}})_*(\chi T)$ on that ray is a number
$t_b(F)\geq0$.

Apply the same construction to the integration current and denote its
coefficient on the isolated primitive ray by
\begin{equation}
                              m_b(F)>0.
\label{eq:p213-actual-calibration-coefficient-revised}
\end{equation}
Strict positivity follows directly from the positive local degree in the
canonical calibration of
\cite[Proposition~14.3.1, \S14.3.2, and Definition~14.5.1]{CLD25},
because $F$ is nonconstant on $b$.  We deliberately do not identify
$m_b(F)$ with $q_b(F)$.  Over an inseparable constant-field extension,
the analytic calibration can retain a local-degree factor that disappears
from the normalized slope after completed-perfect-closure reduction.
For example, take $K=\mathbf F_p((t))$, $L=K(t^{1/p})$, and view
$\mathbf A^{1,\mathrm{an}}_L$ over $K$.  On a type--$3$ segment in
$\mathbf G_{m,L}^{\mathrm{an}}$, the coordinate moment has primitive
normalized slope $q_b(F)=1$, whereas the finite map to
$\mathbf A^{1,\mathrm{an}}_K$ has degree $p$ and hence
$m_b(F)=p$ by \cite[Remark~14.4.3]{CLD25}.  The quotient below retains
this factor.

The ratio
\begin{equation}
                              a_b:=t_b(F)/m_b(F)
\label{eq:p213-intrinsic-branch-ratio-revised}
\end{equation}
is independent of $F$ and of the cutoff.  Indeed, concatenate two moments
and, if necessary, append one separator block; use a common source cutoff
equal to one near $y$.  Write $H$ for this common refinement and let $A$
be the integral projection from its target to the target of $F$.  Since
$A n_b(H)=n_b(F)$, the primitive tangents satisfy
\begin{equation}
       A u_b(H)=r u_b(F),\qquad
       r=q_b(F)/q_b(H)\in\mathbf Z_{>0},\qquad
       A_*[u_b(H)]=r[u_b(F)].
\label{eq:p213-refinement-expansion-revised}
\end{equation}
The second identity follows by pulling back a $(1,1)$-form and changing
the integral edge parameter.  Functoriality of pushforward
\cite[\S7.1.11]{CLD25}, together with separation of the source ray, gives
\begin{equation}
                 t_b(F)=r\,t_b(H),\qquad
                 m_b(F)=r\,m_b(H).
\label{eq:p213-refinement-both-currents-revised}
\end{equation}
Thus the quotient in
\eqref{eq:p213-intrinsic-branch-ratio-revised} is unchanged.  Notice that
every inseparable or nilpotent local-degree factor is contained in both
occurrences of $m_b$ and is never discarded.
Replacing a cutoff by a smaller common cutoff does not change the ray-germ
coefficient.  Indeed, the difference of two such cutoffs vanishes near
$y$.  On the finitely many active arms of the common source star, the
image of its compact support is therefore bounded away from the target
vertex; off the variation graph every positive-degree target form pulls
back to zero.  The pushforward of the difference consequently vanishes on
a sufficiently small target germ.  This argument applies to both $T$ and
$\delta_Y$ and completes the proof of intrinsicity.

At a type--$3$ point, the two integral slope vectors are opposite by
\eqref{eq:p213-normalized-node-relation}.  Their greatest common divisors
are equal and their primitive tangents are opposite.  Applying
\eqref{eq:p213-euclidean-vertex-balance-revised} to the cutoff pushforward
first of $T$ and then of $\delta_Y$ gives equality, respectively, of the
two coefficients $t_b$ and of the two coefficients $m_b$.  Hence the two
branch ratios agree.

Let now $y$ be of type~$2$, put $C=C_y$, and retain the positive integers
$d_P$ from \eqref{eq:p213-variation-integral-branch-weight}.  For each
closed point $P\in C$, let $a_P$ be the ratio of its branch.  We also
attach an intrinsic positive calibration coefficient $m_P$ to that
branch.  Choose any separated moment that contains a norm-one lift of a
residual uniformizer at $P$.  Every coordinate has normalized slope
divisible by $d_P$, while the uniformizer coordinate has slope $d_P$.
Consequently the gcd of the moment's slope vector is $d_P$.  If two such
moments are replaced by their common refinement, its gcd is still $d_P$;
hence the expansion integer $r$ in
\eqref{eq:p213-refinement-expansion-revised} is one for both projections.
Equation \eqref{eq:p213-refinement-both-currents-revised} and cutoff
independence show that the integration-current coefficient is the same for
all these moments.  Denote this coefficient by $m_P>0$.  The corresponding
coefficient for $T$ is $a_Pm_P$.

Define two additive functionals on divisors by
\begin{equation}
       L_T([P])=a_Pm_P,
       \qquad
       L_Y([P])=m_P.
\label{eq:p213-divisor-functionals-revised}
\end{equation}
They are nonnegative on effective divisors.  We claim that both vanish on
principal divisors.

Fix $g\in\operatorname{Rat}(C)^\times$.  In one common moment append a
norm-one lift of $g$, a lift of a residual uniformizer at every point in
$\operatorname{Supp}(\operatorname{div}(g))$, and the finite separator
block constructed in the proof of
\cref{lem:p213-variation-graph-source-cutoff}.  On the branch at $P$, all
normalized slopes are divisible by $d_P$ and the uniformizer coordinate
has slope $d_P$.  Thus $q_{b_P}=d_P$, while the $g$-coordinate of the
primitive tangent is $\operatorname{ord}_P(g)$ by
\eqref{eq:p213-variation-residual-normalized-slope}.  Taking the
$g$-coordinate of
\eqref{eq:p213-euclidean-vertex-balance-revised} first for the cutoff
pushforward of $T$ and then for the cutoff pushforward of $\delta_Y$ gives
\begin{equation}
                         L_T(\operatorname{div}(g))=0.
\label{eq:p213-T-principal-zero-revised}
\end{equation}
\begin{equation}
                         L_Y(\operatorname{div}(g))=0.
\label{eq:p213-Y-principal-zero-revised}
\end{equation}

We use the following elementary Riemann--Roch consequence over the actual
constant field of $C$.  If
$L:\operatorname{Div}(C)\to\mathbf R$ is additive, nonnegative on
effective divisors, and zero on principal divisors, then $L=c_L\deg$.
Indeed, if $D$ has degree zero, choose an effective divisor $A$ with
\begin{equation}
                  \deg A+\chi(C,\mathcal O_C)>0.
\label{eq:p213-rr-positive-constant-revised}
\end{equation}
For every $n\in\mathbf Z$, Riemann--Roch gives
\[
 h^0\bigl(C,\mathcal O_C(A+nD)\bigr)
    \geq \deg A+\chi(C,\mathcal O_C)>0.
\]
Thus $A+nD$ is linearly equivalent to an effective divisor, and
$0\leq L(A)+nL(D)$ for every positive and negative integer $n$.
Consequently $L(D)=0$.  The functional factors through the subgroup
$\deg\operatorname{Div}(C)\subseteq\mathbf Z$, and positivity gives
$c_L\geq0$.

Apply this to \eqref{eq:p213-divisor-functionals-revised}, using
\eqref{eq:p213-T-principal-zero-revised} and
\eqref{eq:p213-Y-principal-zero-revised}.  There are constants
$c_T\geq0$ and $c_Y>0$ such that
\[
                  a_Pm_P=c_T\deg(P),
                  \qquad m_P=c_Y\deg(P).
\]
It follows that $a_P=c_T/c_Y$ for every branch at $y$.  This conclusion
uses the actual analytic calibration weights and therefore retains every
generic inseparable or nilpotent multiplicity.

We now identify the current locally.  Let $\alpha$ be an ordinary smooth
$(1,1)$-form with sufficiently small compact support near
$y\in Y^{[2,3]}$, and present it by a finite moment $F_0$ as in
\cite[Proposition~15.1.1(b)]{CLD25}.  Apply
\cref{lem:p213-variation-graph-source-cutoff}, taking the cutoff equal to
one near $\operatorname{supp}(\alpha)$ and shrinking the support of
$\alpha$, if necessary, into the inverse image of the target neighborhood
$O$.  In the local moment presentation, multiply the target representative
by a compactly supported target cutoff which is one on the compact image
of $\operatorname{supp}(\alpha)$.  This does not change its pullback.
On every branch on which $F_0$ varies, the preceding type--$2$ or
type--$3$ argument gives one common ratio $c_y$.  Edges activated only by
the separator block have zero $F_0$-tangent, so the target representative
vanishes on their primitive tangent.  The Euclidean lemma excludes an
additional vertex current.  The defining identity for compactly supported
pushforward therefore gives
\begin{equation}
                         \langle T,\alpha\rangle
                    =c_y\langle\delta_Y,\alpha\rangle.
\label{eq:p213-local-current-equality-revised}
\end{equation}

The locus $Y^{[2,3]}$ is nonempty and path-connected.  Indeed, a connected
normal analytic curve is an arcwise-connected graph by
\cite[Theorem~3.5.1 and \S1.3.5]{Ducros24CurveStructure}; it is not a
one-point graph because it is purely one-dimensional, and every interior
point of an open graph interval is of type $2$ or $3$.  Convexity of
$Y^{[2,3]}$ follows from \cite[\S4.5.11]{Ducros24CurveStructure}.
We spell out the propagation along an edge.  Fix a branch $b$ at a point
$y\in Y^{[2,3]}$.  The residual-uniformizer construction at type~$2$ and
the non-torsion-modulus construction at type~$3$ give an analytic unit
whose modulus varies on $b$.  After appending a separator moment and
shrinking, its variation graph on that arm has no vertex away from $y$,
and the tropical moment is injective and integral affine there.  Choose a
compact subsegment $J$ of the arm and a source cutoff that is one on a
neighborhood of $J$ and has compact support in a slightly longer segment
star.  The argument of
\cref{lem:p213-variation-graph-source-cutoff} shows that the pushforwards
of both $\chi T$ and $\chi\delta_Y$ are closed on a target neighborhood of
the image of $J$.  By
\cref{lem:p213-euclidean-one-dimensional-current-revised}, their
coefficients $t$ and $m$ are separately constant on this connected target
edge.  Their quotient $t/m$ is therefore constant along $J$.  At every
point of $J$ this quotient is the intrinsic ratio
\eqref{eq:p213-intrinsic-branch-ratio-revised}, because the same moment is
admissible there and that ratio is independent of the moment and cutoff.

Thus the ratio is locally constant along every open graph edge.  At a
type--$2$ or type--$3$ vertex the preceding star argument makes all
incident ratios equal.  Covering a path in the connected locus
$Y^{[2,3]}$ by finitely many such edge segments and vertex stars shows
that the constants $c_y$ in
\eqref{eq:p213-local-current-equality-revised} form one constant $c$ on
this locus.

Finally, every compactly supported ordinary $(1,1)$-form on the curve has
support in $Y^{[2,3]}$.  This follows from the local presentation
\cite[Proposition~15.1.1(b)]{CLD25} and the dimension argument in the
proof of \cite[Proposition~15.2.1(a)]{CLD25}.  A finite smooth partition
of unity on its compact support and
\eqref{eq:p213-local-current-equality-revised} give
$T=c\delta_Y$ globally.  By \cite[Example~8.6.4]{CLD25}, the integration
current is Bedford--Taylor, and the defining identity is stable under
nonnegative scalar multiplication.
\end{proof}

% Live-ready quarantine source-aware proof.
% Date: 29 August 2026.
% This file is not a live-book input and makes no novelty or priority claim.
% PROMOTION GATE: the proof is bound to the revised frozen arbitrary
% normal-curve classification and awaits its production verdict.

\begin{proof}[Proof of
\cref{prop:p213-componentwise-arbitrary-valuation}]
We give the reductions in full because the nilpotent calculus, generic
multiplicities, classification, and curvature measures remember different
parts of the analytic structure.

\paragraph{Nilreduction and the ambient geometry.}
An analytic seminorm kills every nilpotent.  Hence $X$ and
$X_{\mathrm{red}}$ have the same underlying topological space, the same
analytic points, and the same completed residue fields.  Affinoid charts of
the reduction are quotients of affinoid charts of $X$ by their nilradicals.
It follows directly that goodness, separatedness, strictness, and pure
dimension one pass to $X_{\mathrm{red}}$.  The immersion $\iota$ is finite
and has empty relative boundary by \cite[\S11.4.4]{CLD25}.  The
boundary-composition identity \cite[(11.4.1.3)]{CLD25} therefore shows that
$X_{\mathrm{red}}$ is boundaryless.

Both spaces are good and topologically separated, hence locally
holomorphically separated by \cite[\S15.2.2]{CLD25} and locally
differentially separated by \cite[Corollary~15.2.4]{CLD25}.  Thus the Radon
measure and Borel-multiplication statements below apply by
\cite[Lemma~7.2.4, Proposition~7.2.8, and Example~7.2.12]{CLD25}.

If $A\to A/I$ is an affinoid quotient with $I$ nilpotent, every unit of
$A/I$ lifts to a unit of $A$: if $ab=1+n$ with $n$ nilpotent, the finite
geometric series for $(1+n)^{-1}$ corrects $b$ to an inverse.  Every
analytic seminorm on $A$ kills $I$, so a lifted unit and its reduction have
the same absolute value at every analytic point.  Thus a moment map on a
reduced chart lifts, after shrinking, to a moment map on $X$ with the same
tropicalization, and every moment on $X$ restricts to one on the reduction.
This is the nilpotent-immersion case of
\cite[Proposition~15.1.3(b) and the following form variant]{CLD25}.

The construction of ordinary CLD forms from moment charts gives
\begin{equation}
 \label{eq:p213-componentwise-av-form-isomorphism}
 \iota^*:A_X^{p,q}\xrightarrow{\ \sim\ }
                    A_{X_{\mathrm{red}}}^{p,q}.
\end{equation}
The equality of tropicalizations makes
\eqref{eq:p213-componentwise-av-form-isomorphism} compatible with wedge
product, $J$, $\mathrm d'$, $\mathrm d''$, support, and every positivity
cone.  For the topology, fix a compact test support.  A compact-test datum
is a finite collection of moment charts and Lagerberg representatives.
Restriction and the lifted charts above are cofinal collections of such
data, and on corresponding charts the tropical images, subdivisions,
representatives, and all $C^r$-seminorms coincide.  Hence
\begin{equation}
 \label{eq:p213-componentwise-av-test-homeomorphism}
 \iota^*:A^{p,q}_{X,c}\xrightarrow{\ \sim\ }
                    A^{p,q}_{X_{\mathrm{red}},c}
\end{equation}
is a homeomorphism of locally convex spaces.

Taking continuous duals gives
\begin{equation}
 \label{eq:p213-componentwise-av-current-isomorphism}
 \iota_*:D^{a,b}(X_{\mathrm{red}})
       \xrightarrow{\ \sim\ }D^{a,b}(X),
 \qquad
 \langle\iota_*S,\alpha\rangle
       =\langle S,\iota^*\alpha\rangle.
\end{equation}
The test-form positivity cones agree, so this isomorphism preserves
positivity.  It preserves $\mathrm d'$- and $\mathrm d''$-closedness because
$\iota^*$ commutes with the two differentials.  It also identifies the weak
topologies and multiplication by smooth functions.

If $S$ is positive, its product with a locally bounded Borel function is
defined through the Radon measures attached to positive test forms, as in
\cite[Example~7.2.12]{CLD25}.  Those test forms and measures agree under
\eqref{eq:p213-componentwise-av-form-isomorphism}, so
$\iota_*(fS)=f\iota_*S$.  The same identity holds by linearity when locally
$S$ is a difference of two positive currents.  It is independent of the
decomposition: an equality between two sums of positive currents is an
equality of all attached measures, and remains so after multiplication by
$f$.  This does not define a Borel product for an arbitrary order-zero
current.

Finally, the first stage in the CLD psh filtration is the same on $X$ and
$X_{\mathrm{red}}$, because smoothness and positivity of
$\mathrm d'\mathrm d''u$ agree.  Successor stages take decreasing limits,
the continuous filtration takes local uniform limits, and limit stages
take local unions.  Transfinite induction identifies the psh and continuous
psh sheaves.  The $G$-form assertion follows from the $G$-form variant after
\cite[Proposition~15.1.3]{CLD25}.  This proves
\cref{itm:p213-componentwise-av:calculus}.

\paragraph{Normalization and the deleted locus.}
Let
\[
 \bar\nu:\coprod_i\widetilde X_i\longrightarrow X_{\mathrm{red}}
\]
be the componentwise normalization.  On a good affinoid domain define
\[
 \mathcal Q:=\bar\nu_*\mathcal O_{\coprod_i\widetilde X_i}
                  /\mathcal O_{X_{\mathrm{red}}},
 \qquad
 \mathfrak c:=\operatorname{Ann}_{\mathcal O_{X_{\mathrm{red}}}}
                  (\mathcal Q).
\]
Let $Z_0=\operatorname{Supp}(\mathcal Q)=V(\mathfrak c)$ and let $Z$ be
the union of $Z_0$ with the pairwise intersections of distinct irreducible
components.  Analytic local rings on a good analytic space are excellent by
\cite[\S11.1.10]{CLD25}; normalization is finite and its components
correspond to the irreducible components by \cite[\S11.3.1]{CLD25}.
Thus $\mathcal Q$ and the conductor are coherent.  Their support misses a
dense open subset of every reduced component and is zero-dimensional.
The same dimension argument applies to a proper intersection of two
components.  Irreducible components are locally finite by
\cite[\S11.2.5]{CLD25}; on a relatively compact affinoid neighborhood only
finitely many occur.  Hence $Z$ is closed and locally finite.

Put
\begin{equation}
 \label{eq:p213-componentwise-av-conductor-preimage}
                         Z_i:=\bar\nu_i^{-1}(Z).
\end{equation}
It is a proper zero-dimensional Zariski closed subset of the normal curve
$\widetilde X_i$, hence is locally finite.  We verify that it consists of
rigid points without using a nontrivial ground valuation or a compact
strict representative.  Suppose that $z\in Z_i$ were nonrigid.  For a
normal analytic curve, \cite[\S4.2.18]{Ducros24CurveStructure} gives
\[
                 \mathcal O_{\widetilde X_i,z}=\kappa(z),
\]
a field.  A coherent ideal defining $Z_i$ near $z$ has proper stalk in
this field and therefore has zero stalk.  Coherence then makes that ideal
zero on a neighborhood of $z$, forcing the supposedly zero-dimensional
closed set to contain a one-dimensional neighborhood, a contradiction.
Thus every point of $Z_i$ is rigid.  Every point of $Z$ has a preimage in
some $Z_i$.  A finite morphism induces a finite extension of completed
residue fields; since the upper field is finite over $K$, its subfield at
the image is finite over $K$ as well.  Thus every point of $Z$ is rigid.

The finite normalization remains strict in the sense used by CLD.  Indeed,
at a point above a strict germ, homogeneous degrees coming from the base
are torsion modulo $|K^\times|$ by
\cite[Lemma~11.8.5]{CLD25}.  Integrality for a finite local algebra makes
the degree of every additional homogeneous element torsion over those
base degrees; hence it too is torsion modulo $|K^\times|$.  The same lemma
then proves strictness of the normalization germ.  The maps
$X_i\to X_{\mathrm{red}}$ and $\widetilde X_i\to X_i$ are finite and have
empty relative boundary.  Formula (11.4.1.3) in \cite{CLD25} shows that
$\widetilde X_i$ is boundaryless.  It is good and separated, normal by
construction, and connected because it normalizes one irreducible
component.

Set
\begin{equation}
 \label{eq:p213-componentwise-av-normal-open}
                         U_i:=\widetilde X_i\setminus Z_i.
\end{equation}
This is nonempty because $Z_i$ is a proper closed subset.  It is connected:
at a rigid point of a normal analytic curve the normalization fibre has
cardinality one, so the point is unibranch by
\cite[Proposition~3.5.2]{Ducros24CurveStructure}.  A connected analytic
curve is connected by injective arcs by
\cite[Theorem~3.5.1 and \S1.3.5]{Ducros24CurveStructure}, while an interior
point of such an arc is pluribranch by
\cite[\S1.3.12]{Ducros24CurveStructure}.  An injective arc joining two
points outside $Z_i$ therefore cannot pass through a point of $Z_i$.
Thus $U_i$ is a nonempty connected good separated normal boundaryless
purely one-dimensional analytic space.  No smoothness property has been
used.

We shall also use that $\delta_{U_i}$ is nonzero.  This follows directly
from the support theorem for integration currents, without any variation
graph.  The space $U_i$ is boundaryless, topologically separated, and
purely one-dimensional, so \cite[Theorem~12.11.1]{CLD25} makes its
associated tropical space very rich.  Therefore
\cite[Corollary~7.3.13]{CLD25} gives
\begin{equation}
 \label{eq:p213-componentwise-av-nonzero-integration}
                         \operatorname{Supp}(\delta_{U_i})=U_i.
\end{equation}
Since $U_i$ is nonempty, $\delta_{U_i}\ne0$.

\paragraph{Generic multiplicities.}
By \cite[Definition~13.2.5 and \S13.5.2]{CLD25},
\[
                   [X]=\sum_i m_i[X_i],
 \qquad
                   [X_{\mathrm{red}}]=\sum_i[X_i].
\]
Integration agrees with the current attached to the fundamental cycle by
\cite[Proposition~16.6.3]{CLD25}.  The normalization is finite and has
calibrated degree one over the normal locus.  Its exceptional image is
contained in $Z$.  An ordinary smooth $(1,1)$-test germ vanishes at a rigid
point: by \cite[Proposition~15.1.1(b)]{CLD25} it is pulled back from a
local ordinary moment, and the dimension argument in the proof of
\cite[Proposition~15.2.1(a)]{CLD25} makes that pullback zero wherever
$d_K<1$.  A rigid point has $d_K=0$ by \cite[(11.2.3.1)]{CLD25}.
Consequently the cycle identities give the current identities
\eqref{eq:p213-componentwise-av-fundamental-cycles} on all of $X$.
Local finiteness of the components makes both sums locally finite.

If every $m_i$ is one, the two currents are equal.  Conversely, restrict an
asserted equality to $X_i\setminus Z$ and transport it to $U_i$.  It becomes
$(m_i-1)\delta_{U_i}=0$.  Equation
\eqref{eq:p213-componentwise-av-nonzero-integration} gives $m_i=1$.
Since $m_i$ is the length of the generic local ring along $X_i$, this is
equivalent to generic reducedness along every component.  This proves
\cref{itm:p213-componentwise-av:integration}.

\paragraph{Current classification and uniqueness.}
Let $T$ be a positive $\mathrm d'$- and $\mathrm d''$-closed current of
bidegree $(0,0)$ on $X$.  By
\eqref{eq:p213-componentwise-av-current-isomorphism}, there is a unique
positive closed current $T_{\mathrm{red}}$ on $X_{\mathrm{red}}$ with
$T=\iota_*T_{\mathrm{red}}$.  The open set
$X_{\mathrm{red}}\setminus Z$ is the disjoint union of the normal component
loci $X_i\setminus Z$, and
\[
 \bar\nu_i|_{U_i}:U_i\xrightarrow{\ \sim\ }X_i\setminus Z.
\]
Transport the restriction of $T_{\mathrm{red}}$ on $X_i\setminus Z$ to a
current $T_i$ on $U_i$.  Restriction preserves positivity and closedness by
\cite[\S7.1.6 and \S\S7.7.1--7.7.3]{CLD25}.  The arbitrary-valuation normal
curve classification
\cref{prop:p213-normal-curve-arbitrary-valuation-current-revised} gives a number
$c_i\in\mathbf R_{\geq0}$ such that
\begin{equation}
 \label{eq:p213-componentwise-av-classification-off-Z}
                            T_i=c_i\delta_{U_i}.
\end{equation}
The coefficient is unique by
\eqref{eq:p213-componentwise-av-nonzero-integration}.

Write $\bar\delta_i=(\bar\nu_i)_*\delta_{\widetilde X_i}$ on
$X_{\mathrm{red}}$ and set
\[
                         S:=T_{\mathrm{red}}-\sum_i c_i\bar\delta_i.
\]
The sum is locally finite, and
\eqref{eq:p213-componentwise-av-classification-off-Z} says that $S$
vanishes on $X_{\mathrm{red}}\setminus Z$.  If $\alpha$ is a compactly
supported ordinary smooth $(1,1)$-test form, the ordinary moment and
dimension argument just used gives
\[
 \operatorname{Supp}(\alpha)
    \subseteq\{x:d_K(x)\geq1\}
    \subseteq X_{\mathrm{red}}\setminus Z.
\]
Thus $\langle S,\alpha\rangle=0$, and $S=0$.  Pushing forward by $\iota$
proves \eqref{eq:p213-componentwise-av-classification}.  Its uniqueness
follows after restriction to $X_i\setminus Z$, transport to $U_i$, and
testing against the form in
\eqref{eq:p213-componentwise-av-nonzero-integration}.  This proves
\cref{itm:p213-componentwise-av:classification}.

\paragraph{Componentwise curvature measures.}
Let $W,U,V$ be as in
\cref{itm:p213-componentwise-av:bedford-taylor}, and regard $W$ also as an
open subset of $X_{\mathrm{red}}$.  Put
\[
             \widetilde W_i:=\bar\nu_i^{-1}(W),
\]
and denote the restrictions of $\nu_i$ and $\bar\nu_i$ by
$\nu_{i,W}^X$ and $\nu_{i,W}^{\mathrm{red}}$, respectively.  The good
separated spaces $\widetilde X_i$ are locally differentially separated by
\cite[\S15.2.2 and Corollary~15.2.4]{CLD25}.  Pullback preserves continuous
psh functions by \cite[Proposition~8.1.5]{CLD25}.  Hence
\begin{equation}
 \label{eq:p213-componentwise-av-measures}
 \mu_i^U:=\mathrm d'\mathrm d''[\bar\nu_i^*U],
 \qquad
 \mu_i^V:=\mathrm d'\mathrm d''[\bar\nu_i^*V]
\end{equation}
are positive Radon measures by \cite[Lemma~8.5.5]{CLD25}.  Put
$\sigma_i:=\mu_i^U-\mu_i^V$.

For every continuous function $F$ on $W$, multiply the two identities in
\eqref{eq:p213-componentwise-av-fundamental-cycles} by $F$.  The product
is defined by \cite[Example~7.2.12]{CLD25}, and direct testing gives
\begin{align}
 [F]_X
 &=\sum_i m_i(\nu_{i,W}^X)_*
       \bigl(((\nu_{i,W}^X)^*F)\delta_{\widetilde W_i}\bigr),
 \label{eq:p213-componentwise-av-function-current-X}\\
 [F]_{X_{\mathrm{red}}}
 &=\sum_i(\nu_{i,W}^{\mathrm{red}})_*
       \bigl(((\nu_{i,W}^{\mathrm{red}})^*F)
                              \delta_{\widetilde W_i}\bigr).
 \label{eq:p213-componentwise-av-function-current-red}
\end{align}
Finite proper pushforward commutes with the current differentials by
transposition; for continuous psh functions this is also
\cite[Proposition~8.5.6]{CLD25}.  Differentiating
\eqref{eq:p213-componentwise-av-function-current-X} for $F=U,V$ shows
that the left equality in
\eqref{eq:p213-componentwise-av-curvature-equivalence} is equivalent to
\begin{equation}
 \label{eq:p213-componentwise-av-weighted-sum}
                 \sum_i m_i(\nu_{i,W}^X)_*\sigma_i=0.
\end{equation}
Since $\nu_{i,W}^X=\iota\circ\nu_{i,W}^{\mathrm{red}}$ and $\iota_*$ is
injective by \eqref{eq:p213-componentwise-av-current-isomorphism}, this is
equivalent to
\begin{equation}
 \label{eq:p213-componentwise-av-weighted-sum-red}
       \sum_i m_i(\nu_{i,W}^{\mathrm{red}})_*\sigma_i=0
       \qquad\text{on }X_{\mathrm{red}}.
\end{equation}
Likewise the right equality is equivalent to
\begin{equation}
 \label{eq:p213-componentwise-av-unweighted-sum}
                 \sum_i(\nu_{i,W}^{\mathrm{red}})_*\sigma_i=0.
\end{equation}

Set $Z_W=Z\cap W$ and $Z_{i,W}=Z_i\cap\widetilde W_i$.  Away from $Z_W$,
the component images are disjoint and each
$\nu_{i,W}^{\mathrm{red}}$ is an analytic isomorphism onto its component
image.  Restricting
\eqref{eq:p213-componentwise-av-weighted-sum-red} to one component image,
and using $m_i>0$, gives
\begin{equation}
 \label{eq:p213-componentwise-av-support}
                    \operatorname{Supp}(\sigma_i)\subseteq Z_{i,W}.
\end{equation}
The same conclusion follows from
\eqref{eq:p213-componentwise-av-unweighted-sum}.

Every point of $Z_{i,W}$ is rigid, and its local ring on
$\widetilde X_i$ is normal.  If $K$ is nontrivially valued, apply
\cref{lem:p21p3-componentwise-normal-rigid-zero-atom}; if $K$ is
trivially valued, apply
\cref{lem:p213-trivial-gauss-rigid-zero-atom}.  Applied to each of the two
continuous psh functions in
\eqref{eq:p213-componentwise-av-measures}, the appropriate lemma gives
\[
                         \mu_i^U(\{z\})
                         =\mu_i^V(\{z\})=0
                         \qquad(z\in Z_{i,W}).
\]
Because $Z_{i,W}$ is locally finite, a signed Radon measure supported there
is the locally finite sum of its point masses.  Equation
\eqref{eq:p213-componentwise-av-support} and the displayed zero-atom
equalities therefore imply
\begin{equation}
 \label{eq:p213-componentwise-av-component-equality}
                            \sigma_i=0
                            \qquad\text{for every }i.
\end{equation}
Starting from either the weighted or the unweighted sum yields
\eqref{eq:p213-componentwise-av-component-equality}; conversely that
equality implies both sums.  This proves
\eqref{eq:p213-componentwise-av-curvature-equivalence}.

Finally put $h=U-V$ and use
\eqref{eq:p213-componentwise-av-classification}.  Multiplication by $h$ is
defined because $T$ is positive and $h$ is locally bounded.  On compact
test supports the component sum is finite.  Proper pushforward,
\eqref{eq:p213-componentwise-av-component-equality}, and linearity give
\begin{align*}
 \mathrm d'\mathrm d''(hT|_W)
 &=\sum_i c_i(\nu_{i,W}^X)_*
       \mathrm d'\mathrm d''[\bar\nu_i^*U-\bar\nu_i^*V]\\
 &=\sum_i c_i(\nu_{i,W}^X)_*(\mu_i^U-\mu_i^V)=0.
\end{align*}
This is \eqref{eq:p213-componentwise-av-bt-conclusion} and proves
\cref{itm:p213-componentwise-av:bedford-taylor}.
\end{proof}
